\documentclass[12pt, a4paper]{article}

% ---------------------------------------------------------------------------
% Encodage et langue
% ---------------------------------------------------------------------------
\usepackage[utf8]{inputenc}
\usepackage[T1]{fontenc}
\usepackage[french]{babel}

% ---------------------------------------------------------------------------
% Mise en page
% ---------------------------------------------------------------------------
\usepackage[top=2.5cm, bottom=2.5cm, left=2.8cm, right=2.8cm]{geometry}
\usepackage{parskip}
\usepackage{microtype}
\usepackage{lmodern}
\usepackage{setspace}
\setstretch{1.1}

% ---------------------------------------------------------------------------
% Mathématiques
% ---------------------------------------------------------------------------
\usepackage{amsmath}
\usepackage{amssymb}
\usepackage{mathtools}
\usepackage{bm}          % vecteurs en gras

% ---------------------------------------------------------------------------
% Tableaux et listes
% ---------------------------------------------------------------------------
\usepackage{booktabs}
\usepackage{tabularx}
\usepackage{array}
\usepackage{multirow}
\usepackage{enumitem}
\setlist[itemize]{itemsep=2pt}
\setlist[enumerate]{itemsep=2pt}

% ---------------------------------------------------------------------------
% Code source
% ---------------------------------------------------------------------------
\usepackage{listings}
\usepackage{xcolor}

\definecolor{codebg}{RGB}{248,248,248}
\definecolor{codeframe}{RGB}{210,210,210}
\definecolor{codecomment}{RGB}{100,100,100}
\definecolor{codekeyword}{RGB}{0,80,180}
\definecolor{codestring}{RGB}{160,30,30}
\definecolor{codenumber}{RGB}{100,100,100}

\lstset{
  backgroundcolor=\color{codebg},
  frame=single,
  rulecolor=\color{codeframe},
  basicstyle=\ttfamily\small,
  keywordstyle=\color{codekeyword}\bfseries,
  stringstyle=\color{codestring},
  commentstyle=\color{codecomment}\itshape,
  numbers=none,
  breaklines=true,
  showstringspaces=false,
  extendedchars=true,
  inputencoding=utf8,
  literate=
    {é}{{\'e}}1 {è}{{\`e}}1 {ê}{{\^e}}1 {à}{{\`a}}1 {â}{{\^a}}1
    {î}{{\^i}}1 {ô}{{\^o}}1 {û}{{\^u}}1 {ù}{{\`u}}1 {ç}{{\c{c}}}1
    {É}{{\'E}}1 {È}{{\`E}}1 {Â}{{\^A}}1 {Î}{{\^I}}1
    {←}{{$\leftarrow$}}1 {→}{{$\rightarrow$}}1
    {↑}{{$\uparrow$}}1 {↓}{{$\downarrow$}}1,
}

% ---------------------------------------------------------------------------
% Encadrés
% ---------------------------------------------------------------------------
\usepackage{tcolorbox}
\tcbuselibrary{skins, breakable}

\newtcolorbox{remarque}[1][]{
  colback=blue!5!white, colframe=blue!50!black,
  fonttitle=\bfseries, title=Remarque, #1
}
\newtcolorbox{exemple}[1][]{
  colback=green!5!white, colframe=green!50!black,
  fonttitle=\bfseries, title=Exemple concret, #1
}
\newtcolorbox{attention}[1][]{
  colback=orange!8!white, colframe=orange!70!black,
  fonttitle=\bfseries, title=À retenir, #1
}

% ---------------------------------------------------------------------------
% Hyperliens et PDF
% ---------------------------------------------------------------------------
\usepackage[hidelinks]{hyperref}
\usepackage{cleveref}

% ---------------------------------------------------------------------------
% En-tête et pied de page
% ---------------------------------------------------------------------------
\usepackage{fancyhdr}
\pagestyle{fancy}
\fancyhf{}
\fancyhead[L]{\small\textit{Ramanalyze — Documentation complète}}
\fancyhead[R]{\small\textit{Alexandre Souchaud}}
\fancyfoot[C]{\thepage}
\renewcommand{\headrulewidth}{0.4pt}

% ---------------------------------------------------------------------------
% Commandes personnalisées
% ---------------------------------------------------------------------------
\newcommand{\X}{\mathbf{X}}
\newcommand{\Xc}{\tilde{\mathbf{X}}}
\newcommand{\C}{\mathbf{C}}
\newcommand{\vv}{\mathbf{v}}
\newcommand{\s}{\mathbf{s}}
\newcommand{\xbar}{\bar{\mathbf{x}}}
\newcommand{\cm}{cm$^{-1}$}
\newcommand{\uM}{\,\mu{}M}
\newcommand{\nM}{\,nM}

% ---------------------------------------------------------------------------
% Titre
% ---------------------------------------------------------------------------
\title{%
  \textbf{\LARGE Ramanalyze}\\[0.5em]
  \Large Application d'analyse de spectres Raman\\[0.8em]
  \large Documentation technique et guide utilisateur complet
}
\author{Alexandre Souchaud \\ \small CitizenSers}
\date{Mars 2026}

% ===========================================================================
\begin{document}
% ===========================================================================

\maketitle
\thispagestyle{fancy}

\begin{abstract}
\textbf{Ramanalyze} est une application de bureau développée avec PySide6
(Qt), pandas, NumPy, Plotly et scikit-learn. Elle permet d'assembler, de
visualiser et d'analyser des spectres de spectroscopie Raman et SERS à partir
de fichiers \texttt{.txt} bruts et de métadonnées au format Excel ou CSV.

Ce document constitue la documentation complète du logiciel : guide
d'utilisation pas à pas, description de l'architecture logicielle, et
introduction mathématique détaillée à l'Analyse en Composantes Principales
(PCA) telle qu'elle est implémentée dans l'onglet \textit{Exploration}. La
section PCA est rédigée de façon pédagogique, avec des exemples numériques
concrets, pour permettre à tout utilisateur de comprendre en profondeur ce que
les composantes principales signifient physiquement et chimiquement.
\end{abstract}

\tableofcontents
\newpage

% ===========================================================================
\section{Aperçu général}
% ===========================================================================

\textbf{Ramanalyze} guide l'utilisateur à travers les étapes suivantes :

\begin{enumerate}
  \item \textbf{Sélection} des fichiers spectres \texttt{.txt}.
  \item \textbf{Association} avec un fichier de métadonnées (Excel ou CSV).
  \item \textbf{Assemblage} et correction de ligne de base (\textit{baseline})
        par l'algorithme \textit{modpoly} de la bibliothèque \texttt{pybaselines}.
  \item \textbf{Visualisation} interactive des spectres corrigés (Plotly).
  \item \textbf{Analyse} des pics : intensités maximales locales et ratios entre
        pics.
  \item \textbf{Exploration multivariée} par Analyse en Composantes Principales
        (PCA corrélée et PCA libre).
  \item \textbf{Export} des résultats en Excel ou CSV.
\end{enumerate}

% ===========================================================================
\section{Fonctionnalités principales}
% ===========================================================================

\begin{itemize}
  \item Sélecteur de fichiers \texttt{.txt} avec navigation arborescente,
        ajout/retrait, liste des éléments sélectionnés.
  \item Lecture de métadonnées en \textbf{Excel} (\texttt{.xlsx/.xls}) et
        \textbf{CSV} (séparateur détecté automatiquement parmi \texttt{,},
        \texttt{;}, tabulation).
  \item Correction de ligne de base (algorithme \textit{modpoly}, ordre~5 par
        défaut) via \texttt{pybaselines}.
  \item Assemblage des données (spectres + métadonnées) avec prévisualisation
        immédiate.
  \item Sauvegarde du fichier combiné au choix en CSV ou Excel.
  \item Tracé interactif Plotly : zoom, panoramique, légende cliquable,
        export PNG.
  \item Analyse des pics : intensité maximale dans une fenêtre
        $[\nu_0 \pm \Delta\nu]$, ratios automatiques pour toutes les paires
        de pics, graphique interactif des ratios en fonction d'une variable
        de métadonnées.
  \item Export Excel des résultats (feuilles \texttt{intensites} et
        \texttt{ratios}).
  \item PCA corrélée sur les intensités aux pics sélectionnés.
  \item PCA libre sur les spectres complets (matrice normalisée), avec
        sous-onglets \textit{Scores}, \textit{Loadings}, \textit{Reconstruction}.
  \item Scatter matrix (matrice de corrélation entre les pics).
  \item Génération de volumes de titrant selon une distribution gaussienne
        (outil \textit{GaussianVolumesDialog}).
\end{itemize}

% ===========================================================================
\section{Structure du dépôt}
% ===========================================================================

\begin{lstlisting}
spectroscopie_raman/
|-- main.py             # Point d'entree Qt : fenetre principale, onglets
|-- file_picker.py      # Onglet Fichiers : navigation, selection .txt
|-- metadata_creator.py # Onglet Metadonnees : UI Qt (volumes, correspondance)
|-- metadata_model.py   # Couche donnees pure (sans Qt) : calculs, fusion
|-- spectra_plot.py     # Onglet Spectres : trace Plotly interactif
|-- analysis_tab.py     # Onglet Analyse : pics, ratios, export Excel
|-- peak_selector.py    # Onglet Exploration : PCA correlee + PCA libre
|-- data_processing.py  # Chargement .txt, baseline, assemblage
\end{lstlisting}

Le fichier \texttt{metadata\_model.py} est une \textbf{couche données pure} :
il ne contient aucun import Qt et peut être importé, testé et utilisé
indépendamment de l'interface graphique.

% ===========================================================================
\section{Installation}
% ===========================================================================

\subsection{Prérequis}

\begin{itemize}
  \item Python 3.9 ou supérieur.
  \item Système d'exploitation : macOS, Linux ou Windows.
\end{itemize}

\subsection{Environnement virtuel et dépendances}

\begin{lstlisting}[language=bash]
# Creer un environnement virtuel (recommande)
python -m venv .venv
source .venv/bin/activate      # macOS / Linux
# .venv\Scripts\activate       # Windows (PowerShell ou CMD)

# Installer les dependances
pip install PySide6 pandas numpy plotly pybaselines openpyxl scikit-learn scipy
\end{lstlisting}

\begin{remarque}
\texttt{openpyxl} est nécessaire pour lire et écrire les fichiers Excel.
\texttt{PySide6} peut nécessiter des composants Qt WebEngine supplémentaires
selon le système d'exploitation.
\end{remarque}

\subsection{Lancement}

\begin{lstlisting}[language=bash]
python main.py
\end{lstlisting}

Une fenêtre \textbf{Ramanalyze} s'ouvre avec les onglets : \textit{Fichiers},
\textit{Métadonnées}, \textit{Spectres}, \textit{Analyse}, \textit{Exploration}.

% ===========================================================================
\section{Guide d'utilisation (pas à pas)}
% ===========================================================================

\subsection{Onglet Fichiers}

\begin{enumerate}
  \item Naviguez dans l'arborescence de fichiers (double-clic pour entrer dans
        un dossier).
  \item Sélectionnez un ou plusieurs \textbf{fichiers \texttt{.txt}} de
        spectroscopie Raman.
  \item Cliquez \textbf{Ajouter la sélection} : les fichiers sont placés dans
        la liste de droite.
  \item Utilisez \textbf{Retirer} pour enlever un fichier individuel, ou
        \textbf{Vider la liste} pour recommencer.
  \item Le compteur en bas indique le nombre de fichiers prêts pour l'analyse.
\end{enumerate}

\begin{remarque}
Seuls les fichiers \texttt{.txt} sont listés ; les dossiers restent navigables
mais n'apparaissent pas dans la sélection finale.
\end{remarque}

\subsection{Onglet Métadonnées}

\subsubsection{Saisie des volumes}

Remplissez le \textbf{tableau des volumes} (réactifs $\times$ tubes) :
chaque ligne est un réactif (Solution A, B, C\,\ldots), chaque colonne est
un tube expérimental. Les valeurs sont en \textmu{}L.

Pour générer automatiquement des volumes de titrant distribués autour de
l'équivalence chimique, utilisez le bouton \textbf{Générer volumes (gaussien)}.
Cf.\ section~\ref{sec:gauss} pour le détail mathématique de cet outil.

\subsubsection{Correspondance spectres $\leftrightarrow$ tubes}

Le tableau de correspondance associe chaque fichier spectre à un tube
expérimental. Il peut être rempli manuellement ou généré automatiquement.

\subsubsection{Assemblage}

\begin{enumerate}
  \item Cliquez \textbf{Assembler} :
    \begin{itemize}
      \item lecture de chaque fichier \texttt{.txt},
      \item détection et suppression de l'en-tête (\texttt{Pixel;}),
      \item conversion numérique, tri sur le Raman Shift,
      \item correction de ligne de base (modpoly, ordre~5),
      \item fusion avec les métadonnées sur la clé \texttt{Spectrum name}.
    \end{itemize}
  \item Vérifiez le message de succès (nombre de lignes et colonnes du fichier
        combiné).
  \item Cliquez \textbf{Enregistrer} pour sauvegarder en CSV ou Excel.
\end{enumerate}

\begin{remarque}
Le dernier chemin d'enregistrement est mémorisé pour accélérer les
sauvegardes suivantes.
\end{remarque}

\subsection{Onglet Spectres}

\begin{itemize}
  \item Charge automatiquement le fichier combiné produit à l'étape précédente.
  \item Filtre sur les fichiers sélectionnés dans l'onglet \textit{Fichiers}
        (colonne \texttt{file}).
  \item La légende utilise \texttt{Sample description} si disponible, sinon le
        nom de fichier.
  \item Cliquez \textbf{Tracer avec baseline} pour générer la figure Plotly
        interactive (zoom, pan, export PNG).
\end{itemize}

\subsection{Onglet Analyse}

\begin{enumerate}
  \item Rechargez le fichier combiné si nécessaire.
  \item Choisissez un \textbf{jeu de pics} prédéfini :
    \begin{itemize}
      \item \textbf{532~nm} : 1231, 1327, 1342, 1358, 1450~\cm.
      \item \textbf{785~nm} : 412, 444, 471, 547, 1561~\cm.
    \end{itemize}
  \item Réglez la \textbf{tolérance} $\Delta\nu$ (défaut~: 2{,}00~\cm).
  \item Cliquez \textbf{Analyser les pics} :
    \begin{itemize}
      \item recherche de l'intensité maximale dans $[\nu_0 - \Delta\nu,\,
            \nu_0 + \Delta\nu]$ pour chaque spectre et chaque pic,
      \item calcul de tous les ratios $I_i / I_j$ pour chaque paire de pics
            $(i, j)$,
      \item jointure avec les métadonnées (concentration, tube\,\ldots),
      \item affichage d'un tableau triable et d'un graphique interactif des
            ratios en fonction de \texttt{n(EGTA) (mol)}.
    \end{itemize}
  \item \textbf{Exporter résultats (Excel)\,\ldots} pour sauvegarder :
    \begin{itemize}
      \item feuille \texttt{intensites} : intensité maximale par pic et par
            spectre,
      \item feuille \texttt{ratios} : tous les ratios en format long.
    \end{itemize}
\end{enumerate}

\subsection{Onglet Exploration (PCA)}

L'onglet \textbf{Exploration} regroupe la sélection de pics et deux modes de
PCA, accessibles via des sous-onglets.

\subsubsection{Sous-onglets disponibles}

\begin{table}[h!]
  \centering
  \caption{Sous-onglets de l'onglet Exploration.}
  \label{tab:subtabs}
  \begin{tabularx}{\textwidth}{lX}
    \toprule
    \textbf{Sous-onglet} & \textbf{Description} \\
    \midrule
    Spectres normalisés  & Superposition des spectres après normalisation
                           par leur norme vectorielle \\
    Spectres superposés  & Superposition brute (sans normalisation) \\
    PCA corrélée         & PCA calculée sur la matrice des intensités aux pics
                           sélectionnés uniquement \\
    PCA --- Scores       & Position de chaque spectre dans l'espace défini par
                           les composantes de la PCA libre \\
    PCA --- Loadings     & Poids de chaque nombre d'onde pour chaque composante
                           (affiché comme un spectre) \\
    PCA --- Reconstruction & Reconstruction d'un spectre à partir de $N$
                           composantes seulement \\
    Scatter matrix       & Matrice de nuages de points entre toutes les paires
                           de pics (corrélations visuelles) \\
    \bottomrule
  \end{tabularx}
\end{table}

\subsubsection{Paramètres}

\begin{itemize}
  \item \textbf{Composantes PCA corrélée} : nombre de composantes principales
        pour la PCA sur les intensités aux pics.
  \item \textbf{Composantes PCA libre} : nombre de composantes pour la PCA sur
        les spectres complets (2 à 15).
  \item \textbf{Colorier par} : variable de métadonnées utilisée pour colorer
        les points dans les graphiques de scores, spectres\,\ldots
\end{itemize}

\subsubsection{Workflow typique}

\begin{enumerate}
  \item Rechargez les données.
  \item Définissez vos pics cibles et la tolérance.
  \item Lancez l'analyse.
  \item Consultez le sous-onglet \textbf{PCA --- Scores} (colorié par
        concentration) : les spectres se séparent-ils selon votre variable
        d'intérêt ?
  \item Consultez \textbf{PCA --- Loadings} : à quels nombres d'onde les axes
        sont-ils sensibles ?
  \item Utilisez \textbf{PCA --- Reconstruction} avec $N = 1$, $N = 2$,
        $N = 3$\,\ldots\ pour visualiser ce que chaque composante capture
        réellement.
\end{enumerate}

% ===========================================================================
\section{Comprendre la PCA en spectroscopie Raman}
\label{sec:pca}
% ===========================================================================

\subsection{Motivation : pourquoi la PCA ?}

Un spectre Raman est un vecteur de $p \approx 1\,000$ points, chaque
composante étant l'intensité à un nombre d'onde donné. Si l'on dispose de
$n = 15$ spectres, chaque spectre est un \textbf{point dans un espace à
$p = 1\,000$ dimensions}.

\medskip

Ce que l'on cherche à comprendre expérimentalement :

\begin{itemize}
  \item Les spectres se ressemblent-ils selon la concentration en titrant ?
  \item Y a-t-il des groupes distincts d'échantillons ?
  \item Quelles bandes Raman varient le plus entre les tubes ?
\end{itemize}

Répondre à ces questions directement dans un espace à 1\,000 dimensions est
impossible. La \textbf{PCA} (Analyse en Composantes Principales) résout ce
problème en cherchant les quelques \textbf{directions} dans cet espace qui
capturent le plus de variabilité entre spectres, puis en projetant tous les
spectres sur ces directions pour obtenir une représentation en 2D ou 3D
facilement visualisable.

\begin{attention}
La PCA ne \og perd\fg{} pas l'information arbitrairement : elle choisit les
directions qui \textbf{maximisent la variance conservée}. Les premières
composantes capturent donc l'essentiel de la variabilité réelle.
\end{attention}

\subsection{Calcul pas à pas : de la matrice de spectres aux composantes}

\subsubsection{Étape 1 --- Construction de la matrice de données}

On empile les $n$ spectres en lignes pour former une matrice $\X$ de taille
$n \times p$ :

\[
\X =
\begin{pmatrix}
  x_{1,1} & x_{1,2} & \cdots & x_{1,p} \\
  x_{2,1} & x_{2,2} & \cdots & x_{2,p} \\
  \vdots  & \vdots  & \ddots & \vdots  \\
  x_{n,1} & x_{n,2} & \cdots & x_{n,p}
\end{pmatrix}
\]

où $x_{i,j}$ est l'intensité du spectre $i$ au nombre d'onde $\nu_j$
(après correction de baseline).

\begin{exemple}
Avec $n = 15$ spectres et $p = 1\,000$ nombres d'onde :

\medskip
\begin{center}
\begin{tabular}{lcccc}
\toprule
            & $\nu_1$ & $\nu_2$ & $\cdots$ & $\nu_{1000}$ \\
\midrule
Spectre 1   & 120     & 45      & $\cdots$ & 8   \\
Spectre 2   & 118     & 47      & $\cdots$ & 9   \\
$\vdots$    & $\vdots$& $\vdots$& $\ddots$ & $\vdots$ \\
Spectre 15  & 130     & 41      & $\cdots$ & 7   \\
\bottomrule
\end{tabular}
\end{center}
\end{exemple}

\subsubsection{Étape 2 --- Centrage de la matrice}

Avant tout calcul, on \textbf{centre} la matrice : pour chaque nombre d'onde
$j$, on soustrait la moyenne des $n$ spectres à cette fréquence.

\[
\tilde{x}_{i,j} = x_{i,j} - \bar{x}_j
\quad \text{où} \quad
\bar{x}_j = \frac{1}{n} \sum_{k=1}^{n} x_{k,j}
\]

Après centrage, chaque colonne de $\Xc$ a une \textbf{moyenne nulle}. On
travaille désormais sur les \textbf{écarts} par rapport au spectre moyen
$\xbar = (\bar{x}_1, \ldots, \bar{x}_p)$ : c'est ce qui varie d'un spectre
à l'autre qui est informatif, pas la valeur absolue de l'intensité.

\begin{remarque}
Le centrage est une étape \textbf{obligatoire} en PCA. Sans centrage, la
première composante principale serait simplement la direction du vecteur
moyen $\xbar$ (la valeur absolue la plus grande), ce qui n'a pas de sens
physique.
\end{remarque}

\subsubsection{Étape 3 --- Matrice de covariance}

On calcule la \textbf{matrice de covariance} $\C$ de taille $p \times p$ :

\[
\C = \frac{1}{n-1}\, \Xc^\top \Xc
\]

Chaque élément $C_{i,j}$ est la \textbf{covariance} entre les intensités aux
nombres d'onde $\nu_i$ et $\nu_j$ :

\[
C_{i,j} = \frac{1}{n-1} \sum_{k=1}^{n} \tilde{x}_{k,i}\,\tilde{x}_{k,j}
\]

Interprétation physique :

\begin{itemize}
  \item $C_{i,j} > 0$ : quand l'intensité à $\nu_i$ est élevée (au-dessus de
        sa moyenne), celle à $\nu_j$ l'est aussi. Les deux bandes Raman
        \textbf{varient de concert} — elles sont probablement liées au même
        phénomène chimique.
  \item $C_{i,j} < 0$ : les deux bandes varient en sens opposé.
  \item $C_{i,j} \approx 0$ : les bandes sont indépendantes.
  \item $C_{i,i} = \text{Var}(\nu_i)$ : variance de l'intensité à $\nu_i$
        sur l'ensemble des spectres.
\end{itemize}

La matrice $\C$ encode ainsi \textbf{toutes les corrélations entre bandes
Raman} dans un seul objet mathématique de taille $1\,000 \times 1\,000$.

\subsubsection{Étape 4 --- Décomposition en valeurs propres}

On résout le \textbf{problème aux valeurs propres} de $\C$ :

\[
\C\,\vv_k = \lambda_k\,\vv_k
\]

Cette équation admet $p$ solutions (paires $(\lambda_k, \vv_k)$) :

\begin{itemize}
  \item $\lambda_k \geq 0$ est la \textbf{valeur propre}, proportionnelle à la
        variance expliquée par la $k$-ième composante.
  \item $\vv_k \in \mathbb{R}^p$ est le \textbf{vecteur propre} associé,
        appelé \textbf{loading} de la composante PC$_k$.
        Il est de norme unitaire ($\|\vv_k\| = 1$) et a exactement la même
        longueur que les spectres (un coefficient par nombre d'onde).
\end{itemize}

On trie les paires par valeur propre décroissante :
$\lambda_1 \geq \lambda_2 \geq \cdots \geq \lambda_p \geq 0$.
Le plus grand $\lambda_1$ donne \textbf{PC1}, le second $\lambda_2$ donne
\textbf{PC2}, etc.

La \textbf{fraction de variance expliquée} par la composante $k$ est :

\[
\text{VE}_k = \frac{\lambda_k}{\displaystyle\sum_{j=1}^{p} \lambda_j}
\]

et la variance expliquée \textbf{cumulée} par les $N$ premières composantes :

\[
\text{VE}_{\text{cum}}(N) = \frac{\displaystyle\sum_{k=1}^{N} \lambda_k}
                                  {\displaystyle\sum_{j=1}^{p} \lambda_j}
\]

\begin{exemple}
Si $\text{VE}_1 = 72\,\%$, $\text{VE}_2 = 15\,\%$, $\text{VE}_3 = 6\,\%$,
alors les 3 premières composantes expliquent $72 + 15 + 6 = 93\,\%$ de toute
la variabilité entre spectres. Les 997 composantes restantes partagent les
7\,\% restants (principalement du bruit).
\end{exemple}

\begin{remarque}[title=Pourquoi les valeurs propres de $\C$ ?]
La matrice $\C$ est symétrique positive. Ses vecteurs propres forment une
\textbf{base orthonormée} de $\mathbb{R}^p$ : deux composantes principales
distinctes sont toujours perpendiculaires, ce qui garantit qu'elles capturent
des informations \emph{indépendantes}. C'est une propriété fondamentale de la
PCA : les composantes ne se répètent pas entre elles.
\end{remarque}

\subsubsection{Étape 5 --- Calcul des scores}

Le \textbf{score} du spectre $i$ sur la composante PC$_k$ est le
\textbf{produit scalaire} entre le spectre centré et le loading :

\[
s_{i,k} = \tilde{\mathbf{x}}_i \cdot \vv_k
         = \sum_{j=1}^{p} \tilde{x}_{i,j}\, v_{k,j}
\]

Intuition : on \og projette\fg{} le spectre sur la direction $\vv_k$.
Le résultat est \textbf{un seul nombre} par spectre et par composante.

\medskip

La matrice des scores $\mathbf{S}$ de taille $n \times K$ (avec $K$
composantes retenues) s'écrit :

\[
\mathbf{S} = \Xc \cdot \mathbf{V}
\]

où $\mathbf{V} = [\vv_1 \mid \vv_2 \mid \cdots \mid \vv_K]$ est la matrice
des loadings ($p \times K$).

\begin{exemple}
Avec 15 spectres et 2 composantes retenues :

\medskip
\begin{center}
\begin{tabular}{lrr}
\toprule
             & Score PC1 & Score PC2 \\
\midrule
Spectre 1    & $-3{,}2$  & $0{,}8$   \\
Spectre 2    & $-2{,}8$  & $1{,}1$   \\
Spectre 3    & $-1{,}1$  & $0{,}3$   \\
$\vdots$     & $\vdots$  & $\vdots$  \\
Spectre 15   & $+4{,}1$  & $-0{,}5$  \\
\bottomrule
\end{tabular}
\end{center}

\medskip
On peut maintenant tracer ces 15 spectres dans le plan (PC1, PC2) — en 2
dimensions au lieu de 1\,000.
\end{exemple}

\subsection{Interprétation géométrique}

Imaginons que chaque spectre soit un \textbf{point dans l'espace à $p$
dimensions}. Les $n$ spectres forment un \textbf{nuage de points} dans cet
espace.

Ce nuage n'est pas sphérique : il est allongé dans certaines directions
(là où les spectres varient beaucoup entre eux) et aplati dans d'autres
(là où les spectres se ressemblent).

\begin{itemize}
  \item \textbf{PC1} = l'axe selon lequel le nuage est le \emph{plus allongé}
        (variance maximale).
  \item \textbf{PC2} = l'axe perpendiculaire à PC1 selon lequel le nuage est
        le \emph{deuxième plus allongé}.
  \item \textbf{PC3} = perpendiculaire aux deux précédents, troisième variance
        maximale.
  \item \ldots
\end{itemize}

La PCA trouve la meilleure \og orientation\fg{} pour regarder ce nuage depuis
un espace de faible dimension.

\medskip

\begin{center}
\begin{tabular}{c}
\ttfamily
\begin{tabular}{l}
\hspace{2.5cm}PC2\\
\hspace{2.2cm}$\uparrow$\hspace{0.3cm} $\cdot$\ $\cdot$\ (faible conc.)\\
\hspace{2.2cm}$|$\hspace{0.3cm} $\cdot$\ $\cdot$\ $\cdot$\\
\hspace{2.2cm}$|$\hspace{0.2cm} $\cdot$\ $\cdot$\ $\cdot$\ $\cdot$\\
\hspace{2.2cm}$|$\ $\cdot$\ $\cdot$\ $\cdot$\\
\hspace{2.2cm}$+$-------------------$\rightarrow$ PC1\hspace{0.3cm}(forte conc.)\\
\end{tabular}
\end{tabular}
\end{center}

Chaque point représente un spectre. Sa position sur l'axe PC1 mesure sa
\og ressemblance\fg{} avec le pattern du loading PC1.

\subsection{Ce que signifient PC1, PC2, PC3\,\ldots}

\subsubsection{PC1 : la plus grande source de variabilité}

PC1 capture la direction dans laquelle les spectres varient \textbf{le plus}.
C'est la composante qui explique la plus grande fraction de la variance totale.

\medskip

\textbf{Important :} PC1 n'est pas nécessairement la plus chimiquement
intéressante. Elle capturera simplement la plus grande source de variabilité,
qu'il s'agisse :
\begin{itemize}
  \item d'une \textbf{variation de fluorescence de fond} (ligne de base) entre
        spectres — artefact instrumental courant,
  \item d'une \textbf{variation de puissance laser} entre sessions de mesure,
  \item ou de l'effet chimique recherché (concentration, pH, etc.) si c'est
        bien lui qui domine la variabilité.
\end{itemize}

\subsubsection{PC2 : la deuxième source de variabilité}

Perpendiculaire à PC1 (au sens du produit scalaire, donc mathématiquement
indépendante), PC2 capture la \textbf{deuxième source de variabilité} après
que PC1 a \og absorbé\fg{} la plus grande.

\begin{attention}
Dans une expérience SERS de titration bien menée, si la fluorescence de fond
constitue la principale source de variabilité, \textbf{PC1 capturera cet
artefact} et \textbf{PC2 contiendra l'effet chimique de la concentration en
titrant}. Il faut donc regarder \emph{plusieurs} composantes, pas seulement
PC1.
\end{attention}

\subsubsection{PC3 et suivantes}

Capturent des effets de plus en plus mineurs. En pratique, dans les spectres
Raman/SERS :
\begin{itemize}
  \item PC1--PC2 : les deux principaux effets (ligne de base + chimie, ou deux
        effets chimiques distincts).
  \item PC3--PC5 : effets secondaires, variations locales.
  \item Au-delà de PC5 : \textbf{bruit de mesure} dans la plupart des cas.
\end{itemize}

Le \og coude\fg{} du graphe des valeurs propres (\textit{scree plot}) indique
généralement combien de composantes sont physiquement significatives.

\subsection{Interpréter un loading en détail}

Le loading $\vv_k \in \mathbb{R}^p$ de la composante PC$_k$ est un vecteur
de $p$ coefficients, un par nombre d'onde. On peut le \textbf{tracer en
fonction du Raman Shift} exactement comme un spectre expérimental.

\[
\vv_1 = \bigl[\underbrace{0{,}02}_{\nu_1},\;
              \underbrace{-0{,}01}_{\nu_2},\;
              0{,}00,\;
              \underbrace{0{,}15}_{\nu_4},\;
              \ldots,\;
              \underbrace{0{,}08}_{\nu_{1000}}\bigr]
\]

\paragraph{Règle d'interprétation des coefficients du loading :}

\begin{itemize}
  \item Coefficient \textbf{grand positif} à $\nu_j$ : PC1 est très sensible
        à l'intensité à $\nu_j$. Les spectres avec un \emph{score PC1 positif}
        ont une intensité élevée à cette bande.

  \item Coefficient \textbf{proche de zéro} à $\nu_j$ : PC1 est insensible à
        cette bande. Elle ne contribue pas à différencier les spectres selon
        PC1.

  \item Coefficient \textbf{négatif} à $\nu_j$ : les spectres avec un
        \emph{score PC1 négatif} ont une intensité élevée à $\nu_j$, et
        vice-versa.
\end{itemize}

\paragraph{Méthode pratique :}
\begin{enumerate}
  \item Tracer le loading comme un spectre ($v_{k,j}$ en ordonnée,
        $\nu_j$ en abscisse).
  \item Identifier les pics et les creux, exactement comme pour un spectre
        expérimental.
  \item Associer chaque pic à un mode de vibration connu (bases de données
        Raman, littérature).
  \item \textbf{Conclusion} : la composante PC$_k$ est dominée par les modes
        de vibration correspondant aux pics du loading.
\end{enumerate}

\begin{exemple}
Si le loading PC2 montre des pics à 1\,100 et 1\,340~\cm, et que ces bandes
correspondent à des modes de vibration d'un ligand qui chélate le Cu$^{2+}$,
alors \textbf{PC2 capture la réaction de complexation} entre le ligand et le
Cu$^{2+}$ — i.e., l'effet chimique de la titration.
\end{exemple}

\subsection{Interpréter un score en détail}

Le score $s_{i,k}$ du spectre $i$ sur PC$_k$ est sa \og coordonnée\fg{} sur
l'axe de la $k$-ième composante.

\begin{itemize}
  \item Score \textbf{très positif} : ce spectre ressemble fortement au pattern
        décrit par le loading PC$_k$ (intensités élevées aux nombres d'onde à
        coefficients positifs dans le loading, et basses là où les coefficients
        sont négatifs).

  \item Score \textbf{proche de zéro} : ce spectre est \og neutre\fg{} par
        rapport à PC$_k$. Il ne se distingue pas des autres selon cet axe.

  \item Score \textbf{très négatif} : ce spectre est \og opposé\fg{} au
        pattern de PC$_k$ (intensités élevées là où le loading est négatif).
\end{itemize}

\paragraph{Relation entre score et loading :}
Le score est littéralement le \textbf{produit scalaire} entre le spectre
centré et le loading. C'est une mesure de \textbf{similarité orientée} entre
le spectre et la direction PC$_k$.

\subsection{Graphe Scores PC1 \textit{vs} PC2 : la clé de l'interprétation}

C'est le graphe le plus important de la PCA. On représente chaque spectre
par un point de coordonnées $(s_{i,1}, s_{i,2})$, et on \textbf{colorie} les
points selon une variable de métadonnées (concentration, tube, opérateur\,\ldots).

\paragraph{Questions auxquelles ce graphe répond :}

\begin{enumerate}
  \item \textbf{Les spectres se séparent-ils selon ma variable d'intérêt ?}

    \begin{itemize}
      \item Points séparés \emph{le long de PC1} selon la concentration
            $\Rightarrow$ PC1 capture l'effet chimique de la titration.
      \item Points séparés \emph{le long de PC2} selon la concentration
            $\Rightarrow$ l'effet chimique est la 2\ieme\ source de variance
            (PC1 capturait un artefact).
      \item Points \emph{mélangés aléatoirement} $\Rightarrow$ la variabilité
            principale n'est pas liée à la concentration.
    \end{itemize}

  \item \textbf{Y a-t-il des groupes (clusters) d'échantillons ?}
    Des points regroupés en nuages distincts indiquent des états chimiques
    ou physiques différents.

  \item \textbf{Y a-t-il des spectres aberrants (outliers) ?}
    Un point très éloigné des autres signale un spectre anormal
    (contamination, problème instrumental, mauvaise baseline\,\ldots).
\end{enumerate}

\begin{remarque}
On peut aussi tracer PC1 vs PC3, PC2 vs PC3, etc. Chaque combinaison révèle
une facette différente de la variabilité.
\end{remarque}

\subsection{Reconstruction à partir de $N$ composantes}

L'onglet \textbf{PCA --- Reconstruction} permet de reconstituer un spectre à
partir de $N$ composantes seulement, selon la formule :

\[
\hat{\mathbf{x}}_i = \xbar + \sum_{k=1}^{N} s_{i,k}\,\vv_k
\]

où $\xbar$ est le spectre moyen (ajouté après centrage pour revenir dans
l'espace original).

\paragraph{Comment l'utiliser ?}
\begin{itemize}
  \item \textbf{$N = 1$} : on voit uniquement ce que PC1 a capturé dans ce
        spectre. Si PC1 capturait la ligne de base, le spectre reconstruit
        ressemble à une ligne de base lisse.

  \item \textbf{$N = 2$} : on ajoute la contribution de PC2. Si PC2 capturait
        l'effet de la concentration, le spectre reconstruit montre maintenant
        les pics chimiques en plus.

  \item \textbf{$N = 3$} : on ajoute PC3, etc.

  \item \textbf{$N = p$ (toutes les composantes)} : on retrouve exactement le
        spectre original (reconstruction parfaite).
\end{itemize}

\paragraph{Intérêt pratique :}
En comparant le spectre reconstruit avec $N$ et $N-1$ composantes, on voit
\textbf{exactement ce qu'ajoute la $N$-ième composante}. C'est la méthode la
plus directe pour comprendre ce que chaque composante capture physiquement.

\subsection{Variance expliquée et scree plot}

La \textbf{variance expliquée cumulée} est :

\[
\text{VE}_{\text{cum}}(N) = \sum_{k=1}^{N} \text{VE}_k
= \frac{\displaystyle\sum_{k=1}^{N} \lambda_k}
       {\displaystyle\sum_{j=1}^{p} \lambda_j}
\]

On cherche généralement le plus petit $N$ tel que $\text{VE}_{\text{cum}}(N)
\geq 90\,\%$ (ou 95\,\% selon les usages).

\medskip

Le \textbf{scree plot} est le graphe des valeurs propres $\lambda_k$ en
fonction du rang $k$. On cherche le \og coude\fg{} : le point où la courbe
change nettement de pente, indiquant que les composantes suivantes capturent
surtout du bruit.

\begin{exemple}
\begin{center}
\begin{tabular}{lrr}
\toprule
Composante & VE (\%) & VE cumulée (\%) \\
\midrule
PC1 & 72{,}4 & 72{,}4 \\
PC2 & 14{,}8 & 87{,}2 \\
PC3 &  5{,}6 & 92{,}8 \\
PC4 &  2{,}1 & 94{,}9 \\
PC5 &  1{,}3 & 96{,}2 \\
\midrule
PC6--PC15 &  3{,}8 & 100{,}0 \\
\bottomrule
\end{tabular}
\end{center}

\medskip
Le coude est entre PC3 et PC4 : 3 composantes expliquent 92{,}8\,\% de la
variance, et PC4--PC15 n'apportent que 7{,}2\,\% supplémentaires. On retiendra
donc 3 composantes.
\end{exemple}

\subsection{PCA corrélée \textit{vs} PCA libre}

L'onglet Exploration propose deux modes de PCA :

\subsubsection{PCA corrélée}

La PCA est appliquée uniquement sur la \textbf{matrice des intensités aux pics
sélectionnés}. Si 5 pics sont définis, la matrice est $n \times 5$.

\begin{itemize}
  \item \textbf{Avantages} : très rapide, facile à interpréter (5 variables
        seulement), moins sensible au bruit hors des pics.
  \item \textbf{Inconvénients} : dépend fortement du choix des pics. Si un
        pic pertinent est oublié, l'information est perdue.
\end{itemize}

\subsubsection{PCA libre}

La PCA est appliquée sur la \textbf{matrice complète des spectres normalisés}
($n \times p$, avec $p \approx 1\,000$).

\begin{itemize}
  \item \textbf{Avantages} : ne suppose aucun choix a priori de pics.
        L'algorithme trouve lui-même les régions spectrales les plus
        discriminantes (via les loadings). Idéal en mode exploratoire.
  \item \textbf{Inconvénients} : plus sensible au bruit large bande,
        les loadings sont plus difficiles à interpréter car ils impliquent
        1\,000 coefficients.
\end{itemize}

\begin{attention}
Dans la \textbf{PCA libre}, les spectres sont normalisés par leur norme
vectorielle avant la décomposition. Cela garantit que les différences
d'intensité globale (puissance laser variable) ne dominent pas les composantes.
\end{attention}
\newpage
\subsection{Résumé pratique}

\begin{table}[h!]
  \centering
  \caption{Résumé des éléments à analyser en PCA.}
  \label{tab:pca_summary}
  \begin{tabularx}{\textwidth}{lX}
    \toprule
    \textbf{Élément} & \textbf{Ce que ça dit} \\
    \midrule
    Loading PC$_k$ tracé comme spectre
      & Quelles bandes Raman portent la composante $k$. Les pics du loading
        correspondent aux modes de vibration les plus discriminants. \\
    \addlinespace
    Score PC$_k$ d'un spectre
      & À quel point ce spectre ressemble au pattern décrit par le loading
        de PC$_k$. Valeur positive = forte ressemblance, négative = opposé. \\
    \addlinespace
    Graphe scores PC1 vs PC2, coloré par variable
      & La PCA sépare-t-elle les échantillons selon la chimie ? Permet de
        détecter des groupes, des tendances, des outliers. \\
    \addlinespace
    Variance expliquée (ex. 72\,\%)
      & PC1 résume 72\,\% de toute la variabilité entre spectres. Plus c'est
        élevé, plus la composante est importante. \\
    \addlinespace
    Scree plot
      & Graphe des valeurs propres : le coude indique combien de composantes
        sont significatives (au-delà = bruit). \\
    \addlinespace
    Reconstruction avec $N$ composantes
      & Permet de visualiser exactement ce que les $N$ premières composantes
        ont capturé dans un spectre donné. \\
    \bottomrule
  \end{tabularx}
\end{table}

% ===========================================================================
\section{Génération de volumes gaussiens}
\label{sec:gauss}
% ===========================================================================

L'outil \textbf{Générer volumes (gaussien)} (bouton dans l'onglet Métadonnées)
calcule automatiquement les volumes de titrant à pipeter dans chaque tube pour
réaliser une titration SERS autour de l'équivalence chimique. Cette section
décrit en détail chaque paramètre, le calcul des volumes de chaque solution,
et le comportement en mode compte-goutte.

\subsection{Architecture d'un tube expérimental}

Chaque tube contient exactement \textbf{six composants}. Deux varient d'un
tube à l'autre (A et B), les quatre autres sont identiques dans tous les tubes
(C, D, E, F) :

\begin{table}[h!]
  \centering
  \caption{Composition d'un tube expérimental.}
  \label{tab:solutions}
  \begin{tabularx}{\textwidth}{llXl}
    \toprule
    \textbf{Solution} & \textbf{Rôle} & \textbf{Description} & \textbf{Volume} \\
    \midrule
    A & Tampon         & Complément pour maintenir $V_\text{tot}$ constant & Variable \\
    B & Titrant        & Substance dont on fait varier la quantité & Variable \\
    C & Indicateur     & Volume fixe, identique dans tous les tubes & Fixe \\
    D & PEG            & Volume fixe, identique dans tous les tubes & Fixe \\
    E & Nanoparticules & Volume fixe, identique dans tous les tubes & Fixe \\
    F & Crosslinker    & Volume fixe, identique dans tous les tubes & Fixe \\
    \bottomrule
  \end{tabularx}
\end{table}

De plus, l'échantillon lui-même (Cu$^{2+}$) est présent en volume $V_\text{échant}$
identique dans tous les tubes.

\medskip

Le volume total d'un tube $j$ est donc :
\[
V_\text{tot,\,j} = V_\text{échant} + V_{A,j} + V_{B,j}
                   + V_C + V_D + V_E + V_F
\]

Comme les solutions C, D, E, F et l'échantillon ont le même volume dans tous
les tubes, on définit :
\[
V_\text{fixe} = V_\text{échant} + V_C + V_D + V_E + V_F
\]

et le volume disponible pour A et B dans chaque tube est :
\[
V_{AB} = V_\text{tot} - V_\text{fixe}
\]

La contrainte fondamentale est que \textbf{tous les tubes ont le même volume
total} $V_\text{tot}$, ce qui impose :
\[
V_{A,j} + V_{B,j} = V_{AB} \quad \forall j
\]

\subsection{Paramètres du dialogue}

\subsubsection{Paramètres principaux}

\begin{description}
  \item[\textbf{Nombre de tubes} $N$] Nombre total de tubes (2--200,
        défaut~: 11). Inclut éventuellement le tube à 0~µL et le tube
        dupliqué.

  \item[\textbf{C0 Cu (nM)}] Concentration initiale de Cu$^{2+}$ dans
        l'échantillon pur, en nanomolaires. Sert à calculer la quantité de
        matière de Cu$^{2+}$ par tube :
        \[
          n_\text{Cu} = C_0\,[\text{nM}] \times 10^{-9} \times
                        V_\text{échant}\,[\mu\text{L}] \times 10^{-6}
                      \quad [\text{mol}]
        \]

  \item[\textbf{V échantillon (µL)}] Volume d'échantillon par tube,
        identique pour tous. C'est le volume qui fixe $n_\text{Cu}$
        (et donc $V_\text{eq}$).

  \item[\textbf{C titrant (µM)}] Concentration du stock de titrant
        (Solution B), en micromolaires.

  \item[\textbf{V total par tube (µL)}] Volume final souhaité dans chaque
        tube (défaut~: 3\,090~µL). Sert à calculer $V_{AB}$.
\end{description}

\subsubsection{Solutions fixes (C, D, E, F)}

Pour chaque solution fixe, on saisit :
\begin{itemize}
  \item $V_X$ : le volume (µL) à pipeter dans chaque tube,
  \item $C_X^\text{stock}$ : la concentration du stock dans l'unité choisie
        (M, mM, µM, nM ou \% en masse).
\end{itemize}

\paragraph{Contrainte de quantité de matière fixe.}
Lorsqu'on modifie $V_X$ ou $C_X^\text{stock}$, l'application maintient la
\textbf{quantité de matière} $n_X = C_X^\text{stock} \times V_X$ constante
et recalcule l'autre grandeur automatiquement :
\[
\text{si on modifie } C_X \Rightarrow V_X = \frac{n_X}{C_X}
\qquad
\text{si on modifie } V_X \Rightarrow C_X = \frac{n_X}{V_X}
\]

Cela évite de changer accidentellement la quantité de matière injectée dans
les tubes lorsqu'on ajuste la concentration du stock.

\paragraph{Concentrations finales dans les tubes.}
La concentration finale de la solution $X$ dans le tube $j$ est calculée
par dilution :
\[
C_{X,j}^\text{final} = C_X^\text{stock} \times \frac{V_X}{V_\text{tot,\,j}}
\]

\subsubsection{Contraintes de pipetage}

\begin{description}
  \item[\textbf{Marge (µL)}] Volume minimal et maximal pour le titrant. Le
        générateur garantit :
        \[
          V_\text{min} = \Delta V_\text{marge}
          \qquad
          V_\text{max} = V_{AB} - \Delta V_\text{marge}
        \]
        Défaut~: 20~µL. La marge évite les volumes trop petits ou trop
        grands, non pipetables avec précision.

  \item[\textbf{Pas pipette (µL)}] Résolution minimale de la pipette. Tous
        les volumes générés sont arrondis à ce pas (défaut~: 1~µL).

  \item[\textbf{Inclure un tube à 0~µL de titrant}] Si coché, le premier
        tube contient $V_B = 0$ (contrôle sans titrant). Ce tube compte dans
        $N$ : il reste $N - 1$ tubes libres (ou $N - 2$ si \og dupliquer\fg{}
        est aussi coché).

  \item[\textbf{Dupliquer le dernier volume}] Si coché, le dernier tube est
        un réplicat du tube précédant (même $V_B$). Ce tube est annoté
        \og Contrôle\fg{} dans la correspondance spectres~$\leftrightarrow$~tubes.
\end{description}

\subsection{Calcul du volume d'équivalence chimique}

L'équivalence chimique est atteinte lorsque la quantité de matière de titrant
ajoutée égale exactement celle de l'analyte (Cu$^{2+}$). Puisque tous les
tubes contiennent le même volume d'échantillon :
\[
n_\text{Cu} = C_0 \times V_\text{échant}
\]

le volume d'équivalence pour le titrant est :
\[
\boxed{
V_\text{eq} = \frac{n_\text{Cu}}{C_\text{titrant}}
            = \frac{C_0\,[\text{nM}] \times 10^{-9}
                   \times V_\text{échant}\,[\mu\text{L}] \times 10^{-6}}
                  {C_\text{titrant}\,[\mu\text{M}] \times 10^{-6}}
            \times 10^6
            \quad [\mu\text{L}]
}
\]

\begin{exemple}
$C_0 = 1000$~nM, $V_\text{échant} = 1000$~µL, $C_\text{titrant} = 5$~µM~:
\[
n_\text{Cu} = 1000 \times 10^{-9} \times 1000 \times 10^{-6}
            = 10^{-6}~\text{mol}
\]
\[
V_\text{eq} = \frac{10^{-6}}{5 \times 10^{-6}} \times 10^6
            = 200~\mu\text{L}
\]
Le tube d'équivalence reçoit 200~µL de titrant.
\end{exemple}

\subsection{Génération des volumes par quantiles gaussiens}

\subsubsection{Nombre de tubes libres}

Le nombre de volumes \emph{librement calculés} par l'algorithme (ni le
tube à 0~µL ni le tube dupliqué) est :
\[
N_\text{libre} = N - \mathbf{1}[\text{include\_zero}]
                   - \mathbf{1}[\text{duplicate\_last}]
\]

\subsubsection{Quantiles de la loi normale}

On génère $N_\text{libre}$ probabilités uniformément espacées dans $]0,1[$ :
\[
p_i = \frac{i - 0{,}5}{N_\text{libre}}, \quad i = 1, \ldots, N_\text{libre}
\]

puis on les transforme en positions gaussiennes normalisées via la
\textbf{fonction quantile} (inverse de la CDF) de la loi normale standard :
\[
z_i = \Phi^{-1}(p_i)
\]

Exemples de valeurs : $p = 0{,}5 \Rightarrow z = 0$~;
$p = 0{,}16 \Rightarrow z \approx -1$~;
$p = 0{,}84 \Rightarrow z \approx +1$.

Cette transformation produit automatiquement \textbf{plus de points
proches de $z = 0$} (correspondant à $V_\text{eq}$) et moins aux
extrémités, reproduisant le comportement souhaité d'une titration réelle.

\subsubsection{Déformation pour concentrer les points autour de $V_\text{eq}$}

Les quantiles $z_i$ sont déformés par une puissance $\gamma > 0$ :
\[
z'_i = \text{sign}(z_i) \times |z_i|^\gamma
\]

\begin{itemize}
  \item $\gamma = 1$ : distribution gaussienne standard (pas de déformation).
  \item $\gamma > 1$ : les points proches de 0 se resserrent encore davantage
        autour de $V_\text{eq}$. Valeur typique~: $\gamma = 1{,}8$.
  \item $\gamma < 1$ : les points s'espacent au centre.
\end{itemize}

\subsubsection{Facteur d'échelle $\sigma_V$ et bornes physiques}

Le facteur $\sigma_V$ est choisi de façon à ce qu'aucun volume ne sorte
de la plage autorisée $[V_\text{min}, V_\text{max}]$ :
\[
\sigma_V = \xi \times \min\!\left(
  \frac{V_\text{max} - V_\text{eq}}{z'_\text{max}},\;
  \frac{V_\text{eq} - V_\text{min}}{z'_\text{max}}
\right)
\quad \text{avec} \quad z'_\text{max} = \max_i |z'_i|
\]

Le paramètre $\xi \in ]0,1]$ (\texttt{span\_frac}) contrôle le resserrement
global~: $\xi = 1$ fait atteindre les bornes, $\xi < 1$ reserre la
distribution.

\paragraph{Pourquoi le minimum et non le maximum ?}
Utiliser le minimum garantit que les volumes ne \emph{dépassent ni d'un
côté ni de l'autre}. Si $V_\text{eq}$ est proche d'une borne (par exemple
$V_\text{eq} \ll V_{AB}/2$), le minimum contraint $\sigma_V$ du côté le
plus serré, évitant tout dépassement.

\subsubsection{Conversion en volumes finaux}

\paragraph{Volumes de titrant (Solution B) :}
\[
V_{B,i} = V_\text{eq} + \sigma_V \times z'_i
\]

\paragraph{Clip dans les bornes physiques :}
\[
V_{B,i} \leftarrow \text{clip}(V_{B,i},\; V_\text{min},\; V_\text{max})
\]

\paragraph{Quantification sur la grille de pipette :}
\[
V_{B,i}^* = \text{round}\!\left(\frac{V_{B,i}}{\Delta V_\text{pipette}}\right)
             \times \Delta V_\text{pipette}
\]

\paragraph{Ajout du tube à 0~µL et du tube dupliqué (si demandés) :}
\begin{itemize}
  \item Si \texttt{include\_zero}~: $V_{B,0} = 0$ est inséré en premier.
  \item Si \texttt{duplicate\_last}~: $V_{B,N} = V_{B,N-1}$ est ajouté en
        dernier. Ce tube porte le label \og Contrôle\fg{} dans la correspondance.
\end{itemize}

La série finale est triée par ordre croissant.

\paragraph{Volumes de tampon A (Solution A) :}
\[
V_{A,i} = V_{AB}^* - V_{B,i}^*
\]

où $V_{AB}^* = \text{round}(V_{AB} / \Delta V_\text{pipette}) \times
\Delta V_\text{pipette}$ est ajusté sur la grille pour garantir que
$V_{A,i} + V_{B,i}^*$ est identique dans tous les tubes.

\paragraph{Volume total réel :}
\[
V_\text{tot}^\text{réel} = V_\text{fixe} + V_{AB}^*
\]

Ce volume peut légèrement différer de $V_\text{tot}$ demandé (à cause de
l'arrondi sur la grille de pipette).

\subsection{Calcul des concentrations finales dans les tubes}

Pour chaque solution $X$ (C, D, E, F, A, B, échantillon) et chaque tube $j$,
la concentration finale dans la cuvette est :
\[
C_{X,j}^\text{final} = C_X^\text{stock} \times \frac{V_{X,j}}{V_\text{tot,\,j}}
\]

\begin{itemize}
  \item $V_\text{tot,\,j} = V_\text{échant} + V_{A,j} + V_{B,j} + V_C + V_D + V_E + V_F$
  \item Pour les solutions B et A, $V_{X,j}$ varie d'un tube à l'autre.
  \item Pour les solutions C, D, E, F, $V_{X,j} = V_X$ est constant.
  \item $V_\text{tot,\,j}$ est \textbf{identique} dans tous les tubes (grâce
        à la contrainte $V_{A,j} + V_{B,j} = V_{AB}^*$), donc les dilutions
        des solutions fixes sont les mêmes dans tous les tubes.
\end{itemize}

La première ligne du tableau de concentrations indique le volume total de
chaque tube (en µL), les lignes suivantes les concentrations finales (en M)
pour chaque réactif.

\subsection{Mode compte-goutte}

La case à cocher \textbf{Utilisation d'un compte-goutte} reconfigure
automatiquement les paramètres de pipetage pour correspondre à la contrainte
d'un compte-goutte (pas fixe et non ajustable) :

\begin{table}[h!]
  \centering
  \caption{Modifications appliquées par le mode compte-goutte.}
  \label{tab:dropper}
  \begin{tabular}{lll}
    \toprule
    \textbf{Paramètre} & \textbf{Valeur normale} & \textbf{Mode compte-goutte} \\
    \midrule
    Pas pipette ($\Delta V_\text{pipette}$) & 1~µL (ajustable) & 40~µL \\
    Marge ($V_\text{min}$)                 & 20~µL            & 0~µL  \\
    Volume Solution C                      & 30~µL            & 40~µL \\
    Volume Solution F                      & 60~µL            & 40~µL \\
    \bottomrule
  \end{tabular}
\end{table}

\paragraph{Pourquoi pas = 40 µL ?}
Un compte-goutte standard délivre environ 40~µL par goutte. Tous les volumes
doivent donc être des \textbf{multiples de 40~µL}. Le générateur arrondit
chaque $V_{B,i}$ au multiple de 40~µL le plus proche.

\paragraph{Pourquoi marge = 0 ?}
Avec un compte-goutte, la valeur minimale pipetable est 0 (on peut ne pas
mettre de goutte). Il n'y a pas de minimum imposé par la précision de la
pipette.

\paragraph{Volumes C et F ajustés à 40 µL.}
Ces solutions sont souvent ajoutées goutte à goutte également. L'application
les aligne sur le pas du compte-goutte pour que le volume fixe
$V_\text{fixe}$ reste aussi un multiple de 40~µL.

\paragraph{Liberté de modification.}
Après avoir coché la case, tous les champs restent éditables~: l'utilisateur
peut ajuster manuellement chaque valeur si son protocole diffère des valeurs
par défaut du mode compte-goutte.

\begin{remarque}
En mode compte-goutte, la résolution des volumes est plus grossière
(multiples de 40~µL). Le volume d'équivalence $V_\text{eq}$ calculé est
rarement un multiple exact de 40~µL~; le tube \og d'équivalence\fg{}
est donc arrondi au multiple le plus proche. Plus le nombre de tubes est
grand, meilleure sera la couverture autour de $V_\text{eq}$.
\end{remarque}

\subsection{Dialogue \og Quantités de matière\fg{}}

Le bouton \textbf{Quantités de matière\ldots} ouvre un sous-dialogue qui
affiche, pour chaque solution, la quantité de matière $n$ en nmol~:
\[
n_X\,[\text{nmol}] = C_X^\text{stock}\,[\text{M}] \times V_X\,[\mu\text{L}]
                   \times 10^3
\]

L'utilisateur peut modifier directement $n_X$~; l'application recalcule
$C_X^\text{stock}$ en conséquence (avec $V_X$ inchangé). Cette vue est utile
lorsque le protocole est défini en quantités de matière plutôt qu'en
concentrations.

\subsection{Exemple complet numérique}

Protocole~: $N = 5$ tubes, $C_0 = 500$~nM, $V_\text{échant} = 1000$~µL,
$C_\text{titrant} = 2$~µM, $V_\text{tot} = 3090$~µL, $V_C = 30$, $V_D = 750$,
$V_E = 750$, $V_F = 60$~µL, marge = 20~µL, pas = 1~µL, 1 tube à 0~µL,
pas de doublon.

\paragraph{Étape 1 — Volumes fixes :}
\[
V_\text{fixe} = 1000 + 30 + 750 + 750 + 60 = 2590~\mu\text{L}
\]
\[
V_{AB} = 3090 - 2590 = 500~\mu\text{L}
\]
\[
V_\text{min} = 20~\mu\text{L}, \qquad V_\text{max} = 480~\mu\text{L}
\]

\paragraph{Étape 2 — Volume d'équivalence :}
\[
n_\text{Cu} = 500 \times 10^{-9} \times 1000 \times 10^{-6} = 5 \times 10^{-7}~\text{mol}
\]
\[
V_\text{eq} = \frac{5 \times 10^{-7}}{2 \times 10^{-6}} \times 10^6 = 250~\mu\text{L}
\]

\paragraph{Étape 3 — Tubes libres :}
\[
N_\text{libre} = 5 - 1 - 0 = 4 \quad \text{(4 tubes à calculer + 1 tube à 0)}
\]

\paragraph{Étape 4 — Quantiles ($\gamma = 1{,}8$) :}

\begin{center}
\begin{tabular}{ccccc}
\toprule
$i$ & $p_i$ & $z_i = \Phi^{-1}(p_i)$ & $z'_i = \text{sign}(z_i)|z_i|^{1.8}$ & $V_{B,i}$~(µL) \\
\midrule
1 & 0,125 & $-1{,}15$ & $-1{,}44$ & 250 $-$ $\sigma \times 1{,}44$ \\
2 & 0,375 & $-0{,}32$ & $-0{,}12$ & 250 $-$ $\sigma \times 0{,}12$ \\
3 & 0,625 & $+0{,}32$ & $+0{,}12$ & 250 $+$ $\sigma \times 0{,}12$ \\
4 & 0,875 & $+1{,}15$ & $+1{,}44$ & 250 $+$ $\sigma \times 1{,}44$ \\
\bottomrule
\end{tabular}
\end{center}

\medskip
$z'_\text{max} = 1{,}44$, $\sigma_V = \min\!\left(\frac{230}{1{,}44}, \frac{230}{1{,}44}\right) = 160$~µL

\paragraph{Étape 5 — Volumes finaux :}

\begin{center}
\begin{tabular}{lrr}
\toprule
Tube & $V_B$~(µL) & $V_A$~(µL) \\
\midrule
Tube 1 (zéro)    &    0 & 500 \\
Tube 2            &   19 & 481 \\
Tube 3            &  231 & 269 \\
Tube 4            &  269 & 231 \\
Tube 5            &  481 &  19 \\
\bottomrule
\end{tabular}
\end{center}

Les tubes 3 et 4 encadrent $V_\text{eq} = 250$~µL de près~; les tubes 2 et 5
couvrent les extrémités. $V_{A,j} + V_{B,j} = 500$~µL dans tous les tubes.

% ===========================================================================
\section{Formats de fichiers attendus}
% ===========================================================================

\subsection{Fichiers spectres \texttt{.txt}}

\begin{itemize}
  \item En-tête contenant une ligne qui commence par \texttt{Pixel;}.
        Toutes les lignes avant cette ligne sont ignorées.
  \item \textbf{Séparateur de colonnes} : point-virgule (\texttt{;}).
  \item \textbf{Séparateur décimal} : virgule (\texttt{,}).
  \item Colonnes requises après lecture :
    \begin{itemize}
      \item \texttt{Raman Shift} : nombres d'onde en cm$^{-1}$.
      \item \texttt{Dark Subtracted \#1} : intensité brute (après soustraction
            du fond obscur).
    \end{itemize}
  \item Le pipeline produit automatiquement :
    \begin{itemize}
      \item \texttt{Intensity\_corrected} $=$ \texttt{Dark Subtracted \#1}
            $-$ baseline(modpoly).
      \item \texttt{file} : nom du fichier \texttt{.txt}.
      \item \texttt{Spectrum name} : nom sans l'extension \texttt{.txt}
            (clé de jointure avec les métadonnées).
    \end{itemize}
\end{itemize}

\paragraph{Validations appliquées automatiquement :}
\begin{itemize}
  \item Vérification de la plage physique : $-200 \leq \nu \leq 10\,000$
        cm$^{-1}$.
  \item Détection et moyennage des doublons de Raman Shift.
  \item Tri par Raman Shift croissant.
  \item Suppression des lignes avec des valeurs manquantes dans les colonnes
        essentielles.
\end{itemize}

\subsection{Fichier de métadonnées}

\begin{itemize}
  \item \textbf{Excel} : \texttt{.xlsx} ou \texttt{.xls} (lu via
        \texttt{openpyxl}).
  \item \textbf{CSV} : \texttt{.csv}, séparateur auto-détecté parmi
        \texttt{,}, \texttt{;} et tabulation, encodage UTF-8 recommandé.
  \item \textbf{Colonne obligatoire} : \texttt{Spectrum name} (clé de
        jointure avec les spectres).
  \item Colonnes utiles :
    \texttt{Sample description}, \texttt{n(EGTA) (mol)},
    colonnes de concentration\,\ldots
\end{itemize}

% ===========================================================================
\section{Détails techniques (modules)}
% ===========================================================================

\subsection{\texttt{data\_processing.py}}

Ce module gère le chargement et l'assemblage des données.

\begin{itemize}
  \item \texttt{load\_spectrum\_file(path, poly\_order=5, apply\_baseline=True)} :
    \begin{itemize}
      \item Détecte et ignore l'en-tête jusqu'à la ligne \texttt{Pixel;}.
      \item Lit le fichier avec séparateur \texttt{;} et décimales \texttt{,}.
      \item Convertit en numérique, supprime les NaN, trie sur Raman Shift.
      \item Vérifie la plage physique ($-200$--$10\,000$ cm$^{-1}$) et signale
            les valeurs suspectes.
      \item Détecte et moyenne les doublons de Raman Shift.
      \item Applique la correction de baseline modpoly (si demandée).
      \item Retourne un DataFrame avec \texttt{Raman Shift},
            \texttt{Dark Subtracted \#1}, \texttt{Intensity\_corrected},
            \texttt{file}.
    \end{itemize}

  \item \texttt{build\_combined\_dataframe\_from\_df(txt\_files, metadata\_df, \ldots)} :
    \begin{itemize}
      \item Charge chaque fichier \texttt{.txt} via
            \texttt{load\_spectrum\_file}.
      \item Concatène les spectres valides, ajoute \texttt{Spectrum name}.
      \item Fusion gauche sur \texttt{Spectrum name} avec les métadonnées.
      \item Retire les contrôles BRB si \texttt{exclude\_brb=True}.
    \end{itemize}

  \item \texttt{load\_combined\_df(parent, file\_picker, metadata\_creator)} :
        wrapper gérant toutes les erreurs via \texttt{QMessageBox} (vérifie
        que les métadonnées et les fichiers sont disponibles, appelle
        \texttt{build\_merged\_metadata} puis \texttt{build\_combined\_dataframe\_from\_ui}).
\end{itemize}

\subsection{\texttt{metadata\_model.py}}

Couche données pure (\textbf{aucun import Qt}), testable sans interface
graphique. Sépare la logique métier de l'interface utilisateur.

\begin{itemize}
  \item \texttt{to\_float(x)} : conversion robuste d'une valeur quelconque
        (str avec virgule, espace insécable, NaN, None\,\ldots) en float ou
        None.
  \item \texttt{conc\_to\_M(val, unit)} : convertit une concentration en
        mol/L selon l'unité (M, mM, µM, nM, \ldots).
  \item \texttt{normalize\_tube\_label(label)} : normalise des variantes
        d'étiquettes de tube (\texttt{A5}, \texttt{tube05}, \texttt{5},
        \ldots) vers la forme canonique \texttt{Tube 5}.
  \item \texttt{norm\_text\_key(s)} : clé de comparaison insensible à la
        casse et aux accents (é, è, ê, ô, à, â, î, û\,\ldots).
  \item \texttt{is\_control\_label(label)} : retourne \texttt{True} si le
        label correspond à \og Contrôle\fg{} ou \og Control\fg{} (tube
        réplicat).
  \item \texttt{get\_tube\_columns(df)} : retourne les colonnes \texttt{Tube N}
        d'un DataFrame, triées par numéro $N$.
  \item \texttt{infer\_n\_tubes\_for\_mapping(df\_comp)} : déduit le nombre
        de tubes à partir du tableau des volumes. Fallback : 10 tubes.
  \item \texttt{tube\_merge\_key\_for\_mapping(label, n\_tubes, df\_comp)} :
        retourne la clé de jointure entre la correspondance spectres $\leftrightarrow$
        tubes et le tableau des volumes (\texttt{Contrôle} $\to$ \texttt{Tube N}).
  \item \texttt{compute\_concentration\_table(df\_comp)} : calcule les
        concentrations finales dans chaque tube à partir des volumes et des
        concentrations stock.
  \item \texttt{build\_merged\_metadata(df\_comp, df\_map)} : fusionne la
        correspondance spectres $\leftrightarrow$ tubes avec les concentrations
        calculées. Détecte automatiquement le titrant.
  \item \texttt{sers\_gaussian\_volumes(\ldots)} : génère la série de volumes
        de titrant selon la méthode des quantiles gaussiens déformés
        (cf.\ section~\ref{sec:gauss}).
\end{itemize}

\subsection{\texttt{metadata\_creator.py}}

Interface Qt pour la saisie des métadonnées. Délègue tous les calculs purs
à \texttt{metadata\_model} (via des méthodes d'une ligne du type
\texttt{return mm.fonction(...)}). Contient :

\begin{itemize}
  \item Tableau des volumes (réactifs $\times$ tubes) avec éditeur interactif
        (\texttt{TableEditorDialog}).
  \item Tableau de correspondance spectres $\leftrightarrow$ tubes, avec
        synchronisation automatique lors de la modification du tableau des
        volumes.
  \item Dialogue de génération de volumes gaussiens
        (\texttt{GaussianVolumesDialog}).
  \item Sauvegarde et chargement des métadonnées au format \texttt{.xlsx}.
  \item Méthode \texttt{build\_merged\_metadata()} exposée publiquement pour
        les autres onglets.
\end{itemize}

\subsection{\texttt{spectra\_plot.py}}

\begin{itemize}
  \item Charge le fichier combiné depuis \texttt{metadata\_creator}.
  \item Filtre sur les fichiers sélectionnés (colonne \texttt{file}).
  \item Génère une figure Plotly interactive : une courbe par spectre,
        légende = \texttt{Sample description} si disponible.
  \item Export PNG intégré dans l'interface Plotly.
\end{itemize}

\subsection{\texttt{analysis\_tab.py}}

\begin{itemize}
  \item Recharge le fichier combiné via \texttt{load\_combined\_df}.
  \item Propose des jeux de pics prédéfinis (532~nm, 785~nm) ou une saisie
        manuelle.
  \item Calcule l'intensité maximale dans la fenêtre $[\nu_0 \pm \Delta\nu]$
        pour chaque spectre et chaque pic.
  \item Calcule tous les ratios $I_i / I_j$ pour chaque paire de pics $(i, j)$.
  \item Fusionne avec les métadonnées pour afficher les résultats en contexte.
  \item Graphique interactif des ratios vs une variable de métadonnées.
  \item Export Excel avec deux feuilles : \texttt{intensites} et
        \texttt{ratios}.
\end{itemize}

\subsection{\texttt{peak\_selector.py}}

\begin{itemize}
  \item Onglet \textbf{Exploration} : normalisation des spectres (norme
        vectorielle), sélection des pics cibles, tolérance réglable.
  \item \textbf{PCA corrélée} : PCA sur la matrice $n \times p_\text{pics}$
        des intensités aux pics.
  \item \textbf{PCA libre} : PCA sur la matrice normalisée complète
        ($n \times p$). Les données traitées ($\X$, loadings, scores) sont
        stockées comme attributs d'instance pour permettre la reconstruction.
  \item Sous-onglets :
    \begin{itemize}
      \item \textit{Scores} : graphe interactif de scores PC$_i$ vs PC$_j$,
            coloré par variable de métadonnées.
      \item \textit{Loadings} : graphe des coefficients du loading pour chaque
            composante, en fonction du Raman Shift.
      \item \textit{Reconstruction} : reconstruction d'un spectre sélectionnable
            à partir de $N$ composantes, avec affichage simultané du spectre
            original.
    \end{itemize}
  \item Scatter matrix (matrice de nuages de points entre pics).
\end{itemize}

% ===========================================================================
\section{Export et sauvegardes}
% ===========================================================================

\begin{itemize}
  \item \textbf{Fichier combiné} : CSV (\texttt{.csv}) ou Excel
        (\texttt{.xlsx}) au choix dans la boîte de dialogue. Le dernier
        chemin est mémorisé.
  \item \textbf{Résultats d'analyse} : toujours Excel (\texttt{.xlsx}),
        deux feuilles :
    \begin{itemize}
      \item \texttt{intensites} : intensité maximale par pic et par spectre.
      \item \texttt{ratios} : tous les ratios en format long
            (spectre, pic\_num, pic\_den, ratio, métadonnées).
    \end{itemize}
\end{itemize}

% ===========================================================================
\section{Bonnes pratiques}
% ===========================================================================

\begin{enumerate}
  \item S'assurer que \texttt{Spectrum name} dans les métadonnées correspond
        exactement aux noms de fichiers \texttt{.txt} (sans extension).
  \item Utiliser des noms d'échantillons explicites dans \texttt{Sample
        description} pour des légendes lisibles.
  \item Travailler en encodage UTF-8 pour les fichiers CSV.
  \item En PCA, commencer par examiner le graphe Scores coloré par la variable
        d'intérêt (concentration, tube\,\ldots), puis les Loadings pour
        comprendre la signification physique des axes.
  \item Regarder plusieurs composantes (PC1, PC2, PC3) : l'effet chimique
        recherché n'est pas toujours sur PC1.
  \item Utiliser la Reconstruction pour vérifier intuitivement ce que chaque
        composante a capturé.
  \item Vérifier que la correction de baseline est satisfaisante dans
        l'onglet Spectres avant toute analyse quantitative.
\end{enumerate}

% ===========================================================================
\section{Dépannage (FAQ)}
% ===========================================================================

\paragraph{L'assemblage échoue avec un fichier de métadonnées.}
Vérifiez que la colonne \texttt{Spectrum name} est présente et correspond
exactement aux noms de fichiers \texttt{.txt} (sans extension \texttt{.txt},
casse exacte). Pour les CSV, vérifiez le séparateur (\texttt{,}, \texttt{;}
ou tabulation) et l'encodage (UTF-8 recommandé).

\paragraph{Mon fichier combiné n'a pas de colonne \texttt{Intensity\_corrected}.}
Les fichiers \texttt{.txt} doivent contenir les colonnes \texttt{Raman Shift}
et \texttt{Dark Subtracted \#1} après la ligne \texttt{Pixel;}. Vérifiez le
séparateur (\texttt{;}) et le séparateur décimal (\texttt{,}).

\paragraph{Le tracé n'affiche rien dans l'onglet Spectres.}
Vérifiez que :
\begin{enumerate}
  \item l'assemblage a réussi et le fichier combiné a bien été produit,
  \item les fichiers sélectionnés dans l'onglet \textit{Fichiers} existent
        dans la colonne \texttt{file} du combiné,
  \item les colonnes \texttt{Raman Shift} et \texttt{Intensity\_corrected}
        sont présentes.
\end{enumerate}

\paragraph{L'analyse des pics signale des colonnes manquantes.}
Le fichier combiné doit au minimum contenir \texttt{Raman Shift},
\texttt{Intensity\_corrected} et \texttt{file}. Pour le graphique des ratios,
la variable de métadonnées sélectionnée (ex. \texttt{n(EGTA) (mol)}) doit
également être présente.

\paragraph{La PCA n'affiche rien ou plante.}
Il faut au moins 2 spectres et que le nombre de composantes demandé ne
dépasse pas le nombre de spectres (ou le nombre de pics pour la PCA
corrélée). Vérifiez que l'assemblage a réussi avant de lancer l'analyse.

\paragraph{Les spectres contrôle n'apparaissent pas dans la PCA.}
Le label \texttt{Contrôle} dans le tableau de correspondance doit être
orthographié exactement ainsi (avec la majuscule initiale). Il est mappé
automatiquement au dernier tube du tableau des volumes.

% ===========================================================================
\section{Crédits et licence}
% ===========================================================================

\begin{itemize}
  \item \textbf{Développement} : Alexandre Souchaud, CitizenSers.
  \item \textbf{Correction de baseline} : bibliothèque \texttt{pybaselines}
        (\url{https://github.com/derb12/pybaselines}).
  \item \textbf{Interface graphique} : PySide6 / Qt.
  \item \textbf{Graphiques interactifs} : Plotly.
  \item \textbf{PCA} : scikit-learn (\texttt{sklearn.decomposition.PCA}).
  \item \textbf{Statistiques} : NumPy, pandas, SciPy.
\end{itemize}

\medskip
Projet académique / interne --- CitizenSers. Tous droits réservés.

\end{document}
